% --- Archon physics formalization source begin ---
% archon:physics
% archon:covers IPhO2026Problems/problem_IPhO_2026_3_B_1.lean
% archon:source-report reports/ipho_2026/problem_IPhO_2026_3_B_1.source.json
% archon:problem-id IPhO_2026_3
% archon:part-id B.1
% archon:previous-part IPhO_2026_3_A_3
% archon:previous-part-policy natural_language_prerequisite_only

\chapter{Physics problem IPhO\_2026\_3\_B\_1}
\label{ch:IPhO2026Problems_problem_IPhO_2026_3_B_1}

\paragraph{Problem source.}
Continue with the paramagnetic torus.  Its equation of state is T*M*V = n*K*H,
its heat capacity at constant M is C\_M = n*lambda/T\textasciicircum{}2, and dU = C\_M*dT.
The volume is fixed and the magnetic work on the material is
dW = mu\_0*V*H*dM.  Work and heat entering the torus are positive.

Current subquestion:
At fixed temperature T, H changes from H\_i to H\_f. Find the heat Q transferred into the torus.

\paragraph{Current subquestion.}
At fixed temperature T, H changes from H\_i to H\_f. Find the heat Q transferred into the torus.

\paragraph{Recorded answer/context.}
Q = -(mu\_0*n*K/(2*T))*(H\_f\textasciicircum{}2 - H\_i\textasciicircum{}2).

\paragraph{Figure/image path.}
/root/proposal\_for\_physic/science-mango/ipho\_2026\_source/image/T3\_page-3.png

\paragraph{Reusable previous-part conclusions.}
\begin{itemize}
\item Source A.3. Question: Subtract the vacuum-core contribution and write the work dW done on the paramagnetic material. Reusable conclusions: dW = mu\_0*V*H*dM. Policy: natural\_language\_prerequisite\_only; do\_not\_import\_Lean\_output
\end{itemize}

\paragraph{Formalization target.}
create a compiling Lean file with sorry bodies at `IPhO2026Problems/problem\_IPhO\_2026\_3\_B\_1.lean`. The Lean declarations must preserve the physical quantities, dimensions or dimensional roles, figure labels, governing-law hypotheses, and final relation expressed by this problem.
Use Mathlib/Physlib names found through LeanExplore where available. If a physics API is missing, introduce faithful local abstractions rather than scalar placeholder aliases.

\begin{definition}[energyInJoules]
  \label{decl:physics:IPhO_2026_3_B_1:energyInJoules}
  \lean{IPhO2026Problems.IPhO2026_3_B_1.energyInJoules}
  The numerical value, in joules, of a dimensionful energy.
\end{definition}
\begin{proof}
  This is the defining declaration; unfolding it gives the stated typed data or relation.
\end{proof}

\begin{definition}[ParamagneticTorus]
  \label{decl:physics:IPhO_2026_3_B_1:ParamagneticTorus}
  \lean{IPhO2026Problems.IPhO2026_3_B_1.ParamagneticTorus}
  The fixed parameters of the paramagnetic torus.
\end{definition}
\begin{proof}
  This is the defining declaration; unfolding it gives the stated typed data or relation.
\end{proof}

\begin{definition}[ParamagneticTorusLaws]
  \label{decl:physics:IPhO_2026_3_B_1:ParamagneticTorusLaws}
  \lean{IPhO2026Problems.IPhO2026_3_B_1.ParamagneticTorusLaws}
  \uses{decl:physics:IPhO_2026_3_B_1:energyInJoules, decl:physics:IPhO_2026_3_B_1:ParamagneticTorus}
  The governing constitutive and thermodynamic laws used for the torus.
\end{definition}
\begin{proof}
  This is the defining declaration; unfolding it gives the stated typed data or relation.
\end{proof}

\begin{theorem}[Physics formalization target]
\label{thm:physics:IPhO_2026_3_B_1:target}
\lean{IPhO2026Problems.IPhO2026_3_B_1.heatTransferredInto_isothermal}
\uses{decl:physics:IPhO_2026_3_B_1:energyInJoules, decl:physics:IPhO_2026_3_B_1:ParamagneticTorus, decl:physics:IPhO_2026_3_B_1:ParamagneticTorusLaws}
The assigned autoformalize agent should translate this physics problem into Lean declarations in the covered file, with theorem and lemma proof bodies written as `by sorry`.
\end{theorem}
\begin{proof}
This is an autoformalization task, not a proof task. Produce statements that can later be proved by the physics prover without weakening the physical model.
\end{proof}
% --- Archon physics formalization source end ---
